%La línea de abajo es para quitar encabezado
%\thispagestyle{plain}

\chapter*{Introducción}
\markboth{Introducción}{Introducción}
\addcontentsline{toc}{chapter}{Introducción}

En los últimos años, especialmente con la masificación de las redes sociales y plataformas digitales, los procesos electorales han dejado de analizarse únicamente a través de encuestas tradicionales para incorporar nuevas fuentes de información que reflejan el discurso mediático en tiempo real. Plataformas como X/Twitter, YouTube, TikTok y los portales periodísticos digitales generan diariamente grandes volúmenes de contenido relacionado con actores políticos, campañas electorales y debates públicos, configurando ecosistemas digitales altamente interconectados donde se construye dinámicamente la percepción mediática de los candidatos presidenciales.

Cada una de estas fuentes posee características estructurales distintas. Mientras la prensa digital mantiene líneas editoriales relativamente formales, las redes sociales favorecen interacciones inmediatas y contenido de rápida propagación. Del mismo modo, los contenidos audiovisuales incorporan elementos discursivos y emocionales que amplían el impacto comunicacional sobre la audiencia. Como consecuencia, el análisis aislado de una única fuente resulta insuficiente para comprender integralmente la dinámica mediática electoral contemporánea.

A pesar de los avances alcanzados en minería de opiniones y procesamiento de lenguaje natural, gran parte de las investigaciones previas continúan enfocándose en modelos tradicionales de análisis de sentimientos o en enfoques limitados a plataformas específicas. Muchos de estos trabajos presentan dificultades para interpretar contexto político complejo, detectar cambios temáticos en tiempo real o relacionar entidades, eventos y narrativas dentro de un mismo entorno analítico. Asimismo, los métodos convencionales basados exclusivamente en encuestas o análisis estadísticos presentan limitaciones relacionadas con cobertura, periodicidad y capacidad de adaptación frente al comportamiento dinámico de los ecosistemas digitales.

Frente a este escenario, los Modelos de Lenguaje de Gran Escala (LLMs) representan una alternativa prometedora debido a su capacidad para comprender relaciones semánticas complejas, interpretar contexto y procesar lenguaje natural con mayor profundidad que los enfoques clásicos de NLP. Adicionalmente, tecnologías recientes como los grafos de conocimiento, las arquitecturas GraphRAG y los pipelines multimodales permiten integrar información estructurada y no estructurada para generar análisis más contextualizados, explicables y escalables.

En respuesta a esta problemática, la presente investigación propone el desarrollo de un sistema de Inteligencia Electoral basado en LLMs, minería de sentimientos y detección dinámica de temas relevantes, orientado al análisis de la percepción mediática proyectada hacia candidatos presidenciales dentro de ecosistemas digitales. La propuesta integra información proveniente de múltiples fuentes —prensa digital, X/Twitter, YouTube y TikTok— con el propósito de capturar distintas dimensiones del discurso político contemporáneo y superar las limitaciones derivadas del análisis de una sola plataforma.

La arquitectura propuesta incorpora un pipeline de procesamiento de datos basado en el enfoque Medallion, permitiendo organizar las etapas de ingesta, limpieza, enriquecimiento semántico y análisis avanzado del corpus electoral. Sobre esta infraestructura se implementan módulos de análisis de sentimientos asistidos por LLMs, mecanismos híbridos de detección dinámica de temas y procesos de deduplicación semántica para consolidar información relevante dentro de un catálogo temático incremental.

Posteriormente, los resultados obtenidos son integrados en un grafo de conocimiento electoral diseñado para modelar relaciones entre candidatos, medios, temas, sentimientos y eventos relevantes. Esta representación estructurada facilita el análisis de patrones de cobertura, coocurrencia temática y comportamiento editorial entre diferentes medios digitales. Complementariamente, se incorpora un sistema GraphRAG soportado mediante herramientas MCP, permitiendo realizar consultas en lenguaje natural sobre el grafo y generar respuestas contextualizadas a partir de información estructurada y textual.

Asimismo, la propuesta incorpora un enfoque innovador al considerar explícitamente las diferencias en la cobertura mediática entre las distintas fuentes digitales, permitiendo analizar la distribución de la cobertura por medio como un indicador del comportamiento editorial electoral. De esta manera, el principal aporte del estudio radica en la construcción de un sistema integral de Inteligencia Electoral capaz de combinar análisis semántico, modelado relacional y procesamiento multi-fuente dentro de un mismo entorno analítico.

Finalmente, el objetivo principal de la investigación consiste en desarrollar un sistema capaz de analizar de manera más precisa y contextualizada la percepción mediática electoral mediante la combinación de técnicas avanzadas de inteligencia artificial, procesamiento semántico y análisis multi-fuente, contribuyendo tanto al campo de la inteligencia electoral como al estudio computacional de la comunicación política digital en contextos latinoamericanos. La estructura del documento se organiza de la siguiente manera.

En el Capítulo I se desarrolla el planteamiento del problema, incluyendo la descripción de la realidad problemática, la formulación de problemas, los objetivos, hipótesis, justificación y delimitaciones de la investigación.

En el Capítulo II se presentan los antecedentes nacionales e internacionales relacionados con minería de sentimientos, detección dinámica de temas, LLMs, inteligencia electoral y grafos de conocimiento. Asimismo, se exponen las bases teóricas y conceptuales que sustentan la investigación.

En el Capítulo III se describe el entorno empresarial y estratégico asociado al sistema propuesto, incluyendo el análisis organizacional, el modelo de negocio, el análisis FODA y el mapa de procesos del contexto de aplicación.

En el Capítulo IV se presenta la metodología de la investigación y la implementación de la solución propuesta. Se describen el diseño metodológico, la operacionalización de variables y las técnicas de recolección de datos. Asimismo, se detalla la arquitectura del sistema basada en un enfoque Medallion, el pipeline de procesamiento de datos, el framework tecnológico utilizado y la metodología de evaluación del sistema. Finalmente, se expone el cronograma de actividades del proyecto.

En el Capítulo V se desarrolla la implementación de la solución propuesta, abordando las etapas de comprensión del negocio, comprensión de los datos, preparación del corpus, modelamiento, evaluación y despliegue del sistema de inteligencia electoral.

Finalmente, en el Capítulo VI se exponen las conclusiones obtenidas a partir de los resultados de la investigación y se plantean recomendaciones orientadas a futuras mejoras, ampliaciones metodológicas y nuevas líneas de investigación relacionadas con inteligencia artificial aplicada al análisis político digital.

La investigación concluye con las referencias bibliográficas utilizadas y los anexos que complementan la información presentada a lo largo del documento.
